\section{Hyper-Connexiones Restringidas por Variedad (mHC)}
\label{sec:annex-mhc}

Este anexo documenta la integración de las \emph{Hyper-Connexiones Restringidas
por Variedad} (\emph{Manifold-Constrained Hyper-Connections}, mHC)
\citep{xie_mhc_2025}, un marco de conexiones residuales que reemplaza el
residual de identidad estándar por un flujo residual de $n$ streams cuya
matriz de mezcla de flujo se restringe al politopo de Birkhoff. El cuerpo
principal difiere a este anexo la formulación matemática, los detalles de
implementación y los resultados reportados; las claves bibliográficas citadas
aquí se resuelven contra \texttt{docs/bibliography/}.

\subsection{Motivación: la Propiedad de Mapeo de Identidad}

La conexión residual estándar \citep{he_deep_2016} propaga señales a través de
la profundidad sin modificación: con $x_l \in \mathbb{R}^{C}$ y una función de
capa $F$,
\[
x_{l+1} = x_l + F(x_l, W_l),
\]
de modo que, recursivamente, $x_L = x_l + \sum_{i=l}^{L-1} F(x_i, W_i)$. El
término directo $x_l$ (el \emph{mapeo de identidad}) permite que la información
fluya de capas superficiales a profundas sin cambios, lo que He et
al.~\citep{he_deep_2016} identificaron como clave para el entrenamiento estable
de redes muy profundas.

Las \emph{Hyper-Connections} (HC) \citep{xie_mhc_2025} ensanchan el flujo
residual por un factor de expansión $n$: el flujo pasa a ser
$x_l \in \mathbb{R}^{n \times C}$ mientras que la función interna $F$ de cada
capa sigue operando en dimensión $C$. Tres mapeos lineales aprendibles
gobiernan el flujo en cada capa:
\begin{itemize}[leftmargin=*]
    \item $H^{\text{pre}}_l \in \mathbb{R}^{1 \times n}$ agrega el flujo de
        dimensión $nC$ a la entrada de capa de dimensión $C$;
    \item $H^{\text{post}}_l \in \mathbb{R}^{1 \times n}$ proyecta la salida de
        la capa de vuelta al flujo;
    \item $H^{\text{res}}_l \in \mathbb{R}^{n \times n}$ mezcla características
        \emph{dentro} del flujo residual.
\end{itemize}
La propagación de una sola capa es
\[
x_{l+1} = H^{\text{res}}_l x_l + \big(H^{\text{post}}_l\big)^{\!\top}
F\big(H^{\text{pre}}_l x_l, W_l\big),
\]
y, recursivamente, la señal profunda-a-superficial está gobernada por el mapeo
compuesto $\prod_{i=L-1}^{l} H^{\text{res}}_{i+1}$. Debido a que
$H^{\text{res}}_l$ no está restringida en HC, este compuesto se desvía de la
identidad, la norma del flujo no se conserva y la señal hacia adelante/hacia
atrás puede amplificarse o atenuarse sin límite. En un modelo MoE de 27B
parámetros los autores reportan magnitudes de ganancia compuesta con picos
cercanos a $3000$ (frente a $1$ para la identidad), correlacionadas con picos
de norma de gradiente y una subida inesperada de pérdida alrededor del paso
12k. HC además multiplica el costo de acceso a memoria por aproximadamente $n$
(el ``muro de memoria'').

\subsection{Formulación}

Sea $x_l \in \mathbb{R}^{n \times C}$ el flujo en la capa $l$. Los coeficientes
se calculan a partir del flujo aplanado
$\tilde{x}_l = \operatorname{vec}(x_l) \in \mathbb{R}^{1 \times nC}$:
\[
\tilde{H}_l = \frac{1}{r_l}\Big(\alpha \odot
\big(\tilde{x}_l\, \varphi_l\big)\Big) + b_l,
\qquad r_l = \frac{\|\tilde{x}_l\|_2}{\sqrt{nC}},
\]
donde $\varphi_l \in \mathbb{R}^{nC \times (n^2+2n)}$ es una proyección lineal
aprendida, $b_l \in \mathbb{R}^{1 \times (n^2+2n)}$ un sesgo aprendido y
$\alpha^{\text{pre}}, \alpha^{\text{post}}, \alpha^{\text{res}} \in \mathbb{R}$
son escalares de compuerta aprendibles inicializados a un valor pequeño ($0.01$
en el paper y en este código base). El escalado por $1/r_l$ es matemáticamente
equivalente a RMSNorm sobre el flujo aplanado, con la escala por dimensión
absorbida en $\varphi_l$. Los tres bloques de coeficientes son entonces:
\[
H^{\text{pre}}_l = \sigma\big(\tilde{H}^{\text{pre}}_l\big),
\qquad
H^{\text{post}}_l = 2\,\sigma\big(\tilde{H}^{\text{post}}_l\big),
\qquad
H^{\text{res}}_l = \mathrm{SK}\big(\exp(\tilde{H}^{\text{res}}_l)\big),
\]
donde $\sigma$ es la sigmoide logística (que impone no negatividad y evita la
cancelación de señal por coeficientes de signo mixto) y
$\mathrm{SK}(\cdot)$ es la proyección Sinkhorn--Knopp sobre el conjunto de
matrices doblemente estocásticas
\[
\mathcal{M}^{\text{res}} = \big\{\, H \in \mathbb{R}^{n \times n} \;:\;
H \ge 0,\;\; H\,\mathbf{1}_n = \mathbf{1}_n,\;\;
\mathbf{1}_n^{\top} H = \mathbf{1}_n^{\top} \,\big\},
\]
es decir, el politopo de Birkhoff (la envolvente convexa de las matrices de
permutación). Partiendo de $M^{(0)} = \exp(\tilde{H}^{\text{res}}_l)$, la
proyección alterna normalizaciones de filas y columnas,
\[
M^{(t)} = T_c\big(T_r(M^{(t-1)})\big),
\]
para $t = 1, \ldots, t_{\max}$ rondas ($t_{\max} = 20$ por defecto tanto en el
paper como en este código base). Finalmente, la entrada de capa y el flujo
actualizado son
\[
F^{\text{pre}}_l = H^{\text{pre}}_l x_l \in \mathbb{R}^{1 \times C},
\qquad
x_{l+1} = H^{\text{res}}_l x_l +
\big(H^{\text{post}}_l\big)^{\!\top} \otimes F\big(F^{\text{pre}}_l, W_l\big).
\]

\subsection{Propiedades de la Restricción de Variedad}

Restringir $H^{\text{res}}_l$ al politopo de Birkhoff restaura el
comportamiento de conservación del residual de identidad y garantiza tres
propiedades \citep{xie_mhc_2025}:
\begin{enumerate}[leftmargin=*]
    \item \textbf{Preservación de norma}: la norma espectral satisface
        $\|H^{\text{res}}_l\|_2 \le 1$ (mapeo no expansivo), mitigando la
        explosión de gradientes.
    \item \textbf{Cerradura composicional}: las matrices doblemente
        estocásticas son cerradas bajo multiplicación, por lo que el compuesto
        $\prod_i H^{\text{res}}_{i}$ sigue siendo doblemente estocástico a
        profundidad arbitraria, preservando la propiedad de mapeo de identidad
        entre \emph{cualesquiera} dos profundidades.
    \item \textbf{Interpretación geométrica}: como el politopo de Birkhoff es
        la envolvente convexa de las matrices de permutación,
        $H^{\text{res}}_l x$ es una combinación convexa de permutaciones de
        características, es decir, un mecanismo monótono de mezcla de
        información entre streams.
\end{enumerate}
Cuando $n = 1$ la condición doblemente estocástica degenera al escalar $1$,
recuperando el mapeo de identidad exacto. Con $t_{\max}$ finito la proyección
es aproximada, y la ganancia compuesta hacia atrás se desvía de $1$ pero
permanece acotada (alrededor de $1.6$ en las ejecuciones de 27B del paper, una
reducción de ~3 órdenes de magnitud frente a los $\approx 3000$ de HC).

\subsection{Implementación en Frankenstein Transformer}

mHC es un reemplazo opcional de la conexión residual estándar, controlado por
el sub-objeto \texttt{model.mhc} (esquema jerárquico) o por las claves planas
\texttt{use\_mhc}, con \texttt{additionalProperties: false} a nivel de
sub-objeto:

\begin{itemize}[leftmargin=*]
    \item \texttt{enabled} (\texttt{use\_mhc}, bool, por defecto
        \texttt{false}): reemplazar el residual estándar por el flujo residual
        de $n$ streams.
    \item \texttt{expansion\_rate} (\texttt{mhc\_expansion\_rate}, int $\ge 1$,
        por defecto $4$): factor de expansión del flujo $n$; el flujo pasa a
        ser $n \times \texttt{hidden\_size}$ mientras los internos de la capa
        permanecen en \texttt{hidden\_size}. $1$ recupera el mapeo de
        identidad.
    \item \texttt{sinkhorn\_iters} (\texttt{mhc\_sinkhorn\_iters}, int $\ge 1$,
        por defecto $20$): rondas de normalización Sinkhorn--Knopp $t_{\max}$.
    \item \texttt{gating\_init} (\texttt{mhc\_gating\_init}, float $> 0$,
        por defecto $0.01$): valor inicial de los escalares de compuerta
        $\alpha$.
    \item \texttt{checkpoint} (\texttt{mhc\_checkpoint}, bool, por defecto
        \texttt{false}): gradient checkpointing en capas mHC, intercambiando
        cómputo por memoria para mitigar el aumento de $\sim n{\times}$ en la
        memoria de activaciones del flujo residual de $n$ streams.
    \item \texttt{full\_prec\_under\_bitnet}
        (\texttt{mhc\_full\_prec\_under\_bitnet}, bool, por defecto
        \texttt{true}): mantener la proyección de coeficientes $\varphi_l$ como
        \texttt{nn.Linear} de precisión completa incluso cuando
        \texttt{use\_bitnet} es \texttt{true}, evitando el ruido de
        cuantización ternaria en los pequeños coeficientes mHC; cuando es
        \texttt{false} (y BitNet está activado), $\varphi_l$ usa
        \texttt{BitLinear}.
\end{itemize}

El módulo vive en \texttt{src/model/mhc.py}:
\begin{itemize}[leftmargin=*]
    \item \texttt{Sinkhorn\-Knopp\-Function} (\texttt{torch.autograd.Function}):
        el pase hacia adelante aplica las normalizaciones alternadas de filas y
        columnas de $\exp(\cdot)$; el pase hacia atrás recomputa la iteración
        completa bajo \texttt{torch.enable\_grad} y diferencia a través de
        ella, obteniendo el producto Jacobiano-vector exacto (la estrategia de
        recomputación en chip descrita en el paper).
    \item \texttt{ManifoldHyperConnections} (\texttt{nn.Module}): contiene
        $\varphi_l$ (\texttt{proj}), $b_l$ (\texttt{bias}) y los tres escalares
        de compuerta \texttt{alpha\_pre}, \texttt{alpha\_post},
        \texttt{alpha\_res}; expone \texttt{fpre} (proyectar el flujo a la
        entrada de capa), \texttt{recombine} (actualizar el flujo con la salida
        de la capa) y \texttt{mappings} (calcular $H^{\text{pre}},
        H^{\text{post}}, H^{\text{res}}$).
\end{itemize}

Cableado en \texttt{src/model/frankenstein\_model.py}:
\begin{itemize}[leftmargin=*]
    \item \texttt{HybridLayer} gana módulos opcionales \texttt{mhc\_attn} y
        \texttt{mhc\_ffn} (un módulo mHC por función de capa); la ruta
        \texttt{\_forward\_dense\_mhc} ejecuta atención y FFN como funciones de
        capa sobre el flujo compartido $(B, S, n, C)$.
    \item \texttt{FrankensteinTransformer} expande el embedding de dimensión
        $C$ a $(B, S, n, C)$ al entrar mediante \texttt{mhc\_in\_proj} y lo
        colapsa de vuelta a $C$ mediante \texttt{mhc\_out\_proj} antes de la
        cabeza.
    \item Los parámetros mHC se enrutan al grupo de parámetros \texttt{other}
        del optimizador.
    \item mHC es \textbf{incompatible con \texttt{use\_mixture\_of\_depths}}:
        el enrutamiento MoD de tokens opera sobre un flujo $C$-dimensional
        único, en conflicto con el flujo residual de $n$ streams; se lanza un
        \texttt{ValueError} si ambos están habilitados.
\end{itemize}

Ejemplo de configuración: \texttt{configs/examples/es\_arch\_mhc\_adamw.yaml}.

\subsection{Resultados Reportados}

El paper evalúa mHC (y HC) con $n = 4$, $t_{\max} = 20$ e inicialización de
$\alpha$ en $0.01$ sobre modelos MoE estilo DeepSeek-V3 (atención MLA, RoPE,
RMSNorm, balance de carga sin pérdida auxiliar) de 3B a 27B parámetros
\citep{xie_mhc_2025}:
\begin{itemize}[leftmargin=*]
    \item \textbf{Estabilidad}: mHC elimina las subidas de pérdida de HC (p.
        ej., cerca del paso 12k en el modelo de 27B) y mantiene acotada la
        ganancia residual compuesta ($\approx 1.6$ frente a picos de HC de
        $\approx 3000$).
    \item \textbf{Calidad}: en el modelo de 27B, mHC mejora la pérdida final de
        entrenamiento en $0.021$ sobre la línea base, supera a la línea base en
        los 8 benchmarks de evaluación (BBH, DROP, GSM8K, HellaSwag, MATH,
        MMLU, PIQA, TriviaQA) y a HC en 7 de 8, con las mayores ganancias en
        tareas de razonamiento ($+2.1$ BBH, $+2.3$ DROP sobre HC).
    \item \textbf{Escalado}: la mejora relativa de pérdida sobre la línea base
        se mantiene a través de escalas de cómputo (3B $\to$ 9B $\to$ 27B) y a
        través de 1T tokens de escalado de datos.
    \item \textbf{Sobrecarga}: con kernels fusionados (TileLang), cálculo de
        coeficientes en precisión mixta, recomputación selectiva y solapamiento
        de comunicación DualPipe, mHC añade solo $6.7\%$ de tiempo adicional de
        entrenamiento con $n = 4$ en el modelo de 27B.
\end{itemize}
